\section{Artifact Index of the Lean 4 Verification}

Section~\pSecVerify reports the machine verification of the algebraic skeleton. This appendix makes that claim inspectable without access to the artifact: it records where the artifact lives, how it is reproduced, what its axiom budget is, which paper item each machine-checked declaration discharges, and under exactly which premises the load-bearing declarations hold.

\subsection{Layout, Reproduction, and Axiom Budget}

The artifact is a self-contained Lean~4 project distributed in the \texttt{lean-proof} directory of the project repository at
\begin{center}
\texttt{https://github.com/aki-xavier/frozen-axis}
\end{center}
under the \texttt{lean-proof} path. It consists of four files: \texttt{FormalProof.lean}, the single development file of 1219 lines containing every definition, lemma, proposition, theorem, the final theorem, the example variants, and the generalization upgrades; \texttt{lakefile.lean}, which pins Mathlib to the matching release; \texttt{lake-manifest.json}, which locks every transitive dependency to an exact revision; and \texttt{lean-toolchain}, which pins the toolchain to \texttt{leanprover/lean4:v4.34.0}. No component outside this directory is required.

Reproduction is a single command from inside the directory:
\begin{center}
\texttt{cd lean-proof \&\& lake build}
\end{center}
The build was executed on the developer's machine at the blob hash recorded in Section~\pSecVerify and terminates with zero errors, zero warnings, and zero occurrences of \texttt{sorry}; the development also contains no \texttt{axiom} declaration of its own. The reproducible form of this claim, stronger than the \texttt{sorry}-count, is the axiom budget: the command \texttt{\#print axioms} applied to the load-bearing declarations returns
\begin{center}
\texttt{[propext, Classical.choice, Quot.sound]}
\end{center}
for each of \texttt{finale}, \texttt{lemma1\_sharp}, and \texttt{tail\_realization}, and, by the same command, for the remaining declarations of the development. These three are Lean's own standard axioms, the ones every Mathlib development inherits; in particular the list contains no \texttt{sorryAx}, so no step of the chain is discharged by an unproved assertion. The artifact is therefore auditable at the level of the proof term, not merely at the level of a reported build outcome.

One scope remark belongs here, because it delimits what the artifact does and does not certify. The verified content is the deductive chain: definitions, the mask criterion, Lemmas 1 to 3, Proposition 4 and Theorem 5, and their composition in \texttt{finale}. The premises of that chain which concern the modeling side rather than the algebra are not certified, because they are not theorems: the satisfaction of Criterion P4 on real tasks appears as an explicit premise parameter, the continuous-side Euler--Lagrange derivation is not formalized, in the sense of the correspondence that turns the variational objective into pointwise form, while its discrete counterpart, the per-position coefficient (\pEqFoc) of an arbitrary even differentiable potential together with the first-order condition it yields, is machine-checked; and network nondegeneracy is a checkable property. Section~\pSecLayers assigns each component its tier, and Table~\pTabLayers is the authoritative statement of that boundary.

\subsection{Index from Paper Items to Machine-Checked Declarations}

Table~\ref{tab:artifact} lists the correspondence. The tier column repeats the assignment of Table~\pTabLayers so that the reader can see, for each paper item, whether its machine-checked form is carrier-independent (Tier A), anchored on a Mathlib structure theorem (Tier B), or outside the machine guarantee (Tier C).

\begin{table}[htbp]
\caption{Index from paper items to Lean 4 declarations. Declaration names are those of \texttt{FormalProof.lean}. Tier A: carrier-independent algebraic core; Tier B: structured general forms anchored on Mathlib; Tier C: modeling components not covered by the machine guarantee.}
\label{tab:artifact}
\centering
\small
\renewcommand{\arraystretch}{0.88}
\begin{tabular}{@{}>{\raggedright\arraybackslash}p{0.26\textwidth}>{\raggedright\arraybackslash}p{0.56\textwidth}>{\centering\arraybackslash}p{0.09\textwidth}@{}}
\hline
\textbf{Paper item} & \textbf{Lean declaration} & \textbf{Tier}\\
\hline
Definition 1 (ill-posedness), Criterion P1 & \texttt{IllPosed}, \texttt{P1} & B\\
\hline
Mask criterion, both directions & \texttt{illPosed\_of\_hole}, \texttt{not\_illPosed\_of\_full} & B\\
\hline
Mask criterion in kernel form & \texttt{illPosed\_iff\_nontrivial\_ker} & B\\
\hline
Definition 2 (context rule) & \texttt{IndexedContextRule}, \texttt{IsScalarFamily}, \texttt{IsFrozenFamily} & A\\
\hline
Definition 3 (Mode A / Mode B) & \texttt{IsModeA}, \texttt{DependsOnlyOn}, \texttt{IsModeB} & A\\
\hline
Definition 4 (Paradigm E / I) & \texttt{ParadigmE}, \texttt{ParadigmI}, \texttt{IsParadigmI} & A\\
\hline
Criterion P2 (heterogeneity) & \texttt{P2}, \texttt{p2\_gives\_nonconst}, \texttt{hetero\_gives\_nonconst} & B\\
\hline
Criterion P3 (distributional openness; mask form) & \texttt{P3}, \texttt{p3\_implies\_p1} & A\\
\hline
Criterion P4, component form (Section~\pSecPfour) & premise parameter \texttt{hGapPositive}; the barrier, bounded-description, and realization conditions and the barrier-to-gap inequality are prose only & C\\
\hline
Lemma 1 & \texttt{lemma1\_sharp}, \texttt{lemma1\_final}, \texttt{foc\_proportionality} & A\\
\hline
Lemma 1, general form, (\pEqFoc) onward & \texttt{general\ttu proportionality}, \texttt{scalar\ttu family\ttu iff\ttu ratio\ttu agrees}, \texttt{frozen\ttu family\ttu two\ttu configs\ttu iff}, \texttt{position\ttu wise\ttu escape\ttu witness}, \texttt{KappaPhi\ttu quadratic}, \texttt{general\ttu first\ttu order\ttu of\ttu local\ttu min} & A\\
\hline
Lemma 2, extrinsic side & \texttt{lemma2\_full}, \texttt{lemma2\_E\_stationary\_infeasible} & A\\
\hline
Lemma 2, endogenous side & \texttt{paradigmI\_unique\_minimizer}, \texttt{paradigmI\_inner\_unique}, \texttt{paradigmI\_L2\_ae} & A and B\\
\hline
Lemma 3, toy form & \texttt{lemma3\_final2} & B\\
\hline
Lemma 3, general measure form & \texttt{lemma3\_general}, \texttt{tail\_realization} & B\\
\hline
Lemma 3, gap mechanism, both demands of Example~\pExGap & \texttt{toy\_positive\_gap}, \texttt{modeB\_positive\_gap}, \texttt{modeA\_loss\_attained}; \texttt{evidence\_gap\_is\_full}, \texttt{modeB\_evidence\_gap}, \texttt{modeA\_evidence\_attained} & A\\
\hline
Proposition 4 (the network rule is Mode B) & \texttt{neural\_isModeB}, \texttt{neural\_not\_modeA} & A\\
\hline
Nondegeneracy lemma & \texttt{nondeg\_not\_dependsOn} & A\\
\hline
Collision construction & \texttt{collision\_of\_masked}, \texttt{collision\_general} & A and B\\
\hline
Theorem 5 (closure) & \texttt{theorem5} & A\\
\hline
Final theorem & \texttt{finale} & A and B\\
\hline
Continuous-side optimality & Euler--Lagrange derivation; not formalized & C\\
\hline
Network nondegeneracy & checkable property; not formalized & C\\
\hline
\end{tabular}
\end{table}

Four entries in Table~\ref{tab:artifact} deserve emphasis because they are the ones a reader is most likely to look for. First, Criterion P4 is listed as Tier C and as a premise parameter rather than as a formalized definition: the verified form of Lemma 3 consumes \texttt{hGapPositive}, the hypothesis that the risk gap exceeds the bound, and does not attempt to formalize either the criterion or its component form of Section~\pSecPfour, in which the barrier, bounded-description, and realization conditions are combined into the risk-gap floor by the barrier-to-gap inequality. This is the same dislocation already stated in Section~\pSecVerify, recorded here in the index so that it cannot be missed. Second, the collision construction and the classification and closure declarations are Tier A: their proofs quantify over arbitrary index types, so the toy carrier plays no role in them and the toy end and the continuous end are both instantiations of one proof. Third, Definition~2's indexed family is listed as Tier~A rather than as Tier~C: the family, its constant and frozen specializations, and the algebra of the demand ratio are machine-checked, so what a reader cannot verify from the artifact is not the indexed reading but the step in which a general pairwise potential becomes the coefficient (\pEqFoc), together with the operator-valued couplings that the general form does not claim.
Fourth, the general form's premise is a derived statement rather than an assumption: the artifact differentiates the objective along a coordinate line, obtains the per-position coefficient of an arbitrary even differentiable potential, and reaches the first-order condition at a local minimum through the one-dimensional Fermat theorem, with the paper's normalization of constant factors recovered as an explicit scalar. The same derivation is carried out for the quadratic objective, which is where the paper's own stationarity condition of Lemma~1 comes from. 
\subsection{The Premises of the Load-Bearing Declaration}

The index names the declarations; what a reader must inspect to judge the strength of the claim is their premise lists.
The load-bearing composed statement is given verbatim, in its logical content, transliterated to ASCII for typesetting
(\texttt{R} for the reals, \texttt{not}, \texttt{forall}, \texttt{>=}, \texttt{!=} for the corresponding symbols); the
Unicode source of it and of every other statement of the development is in the artifact.

\begin{quote}
\small
\begin{verbatim}
theorem finale (sys : System) (hP3 : P3 sys)
    (m : MeasureSetup) (eps : R) (heps : eps > 0) (W : R) (hW : W != 0)
    (hTail : exists d, m.lossR (neuralRule W) d >= m.riskP (neuralRule W))
    (hGapPositive : m.gapP (neuralRule W) > eps)
    (hNonnegMin : m.riskMin (neuralRule W) >= 0) :
    not IsModeA (neuralRule W)
    and not GeneralizationBound m eps (neuralRule W)
\end{verbatim}
\end{quote}

This premise list is the complete inventory of what the machine-checked composition assumes: Criterion P3 as a hypothesis
on the system, the three measure-theoretic premises of Lemma~3, and the nondegeneracy \texttt{W != 0} of the carrier's
parameter. Nothing else enters, and in particular no premise is smuggled in as a definitional unfolding; the gap premise
is assumed directly, and its derivation from the barrier, bounded-description, and realization conditions of
Section~\pSecPfour is prose. The same reading applies to the other declarations named in Table~\ref{tab:artifact}, whose
premise lists are shorter: the shape exclusions of Lemma~1 and Lemma~2 are implications between stationarity conditions,
with no convexity anywhere; the general form of Lemma~1 replaces the residual by the data-term derivative and the demand
ratio; the general form of Lemma~3 reduces the premises to a risk gap exceeding the bound and nonnegativity of the minimum
risk, because upper-tail realization is discharged as the theorem \texttt{tail\_realization} for an arbitrary probability
measure and an arbitrary integrable function; and the differentiation statements derive the first-order condition from the
objective, which is where analysis rather than algebra enters and where the paper's normalization of constant factors
appears as an explicit scalar. This is the concrete sense in which the artifact is inspectable: the reader can check each
premise list against the prose of Section~\pSecFalsify and see exactly where the two agree and where, in the case of
P4 and in the differentiation of the general pairwise potential, the verified form is explicitly weaker.

\subsection{Maintaining the Correspondence}

The index of Table~\ref{tab:artifact} is kept in correspondence with the sources by the repository's own layout: \texttt{README.md} carries the same index in Markdown, and the single-file submission version of this paper is regenerated from \texttt{sections/} rather than edited by hand. A reader who finds a discrepancy between the index and the source should treat the source of \texttt{FormalProof.lean} as authoritative and report it; the index is a map, not a second copy of the development.